% Propietario conservado: /Users/ruben/Documents/New project/output/REVISION_ARGUMENTAL_APERTURAS_CIERRES_20260915/VI/delivery/MANUSCRIPT_VI_ES_EN_COMPLETE/source_es/sections/13_fock_pauli_composicion.tex
\section{Complétion graduée des canaux et principe d’exclusion}
\label{vi:vi:sec:fock-pauli}

La construction d’une particule fournit un espace de canaux, ses
étiquettes et les opérateurs qui transportent orientation et mémoire. La
composition de plusieurs particules ajoute une question distincte : comment
ces canaux s’échangent et comment leurs occupations sont dénombrées. Ici,
on construit cette strate à partir des données déjà produites par
APP--TRIT--TPK. La masse et l’hamiltonien peuvent agir ensuite sur
l’espace obtenu ; leurs valeurs n’interviennent pas dans la démonstration des
règles d’occupation.

La parité employée est celle du transport de feuillet dans la représentation
d’une particule. Ce n’est ni le signe de la charge, ni la parité d’un entier
visible, ni la troisième coordonnée du TRIT considérée isolément.
Les différentes lectures conservent leurs applications. La
section~\ref{vi:vi:sec:conjugacion} construit le retour orienté et la
représentation spinorielle ; cette section développe l’algèbre
multiparticulaire correspondant à une graduation et à une règle
d’échange déclarées.

\subsection{Espace fini de canaux et caractère de parité}
\label{vi:vi:subsec:canales-graduados}

Une profondeur \(K\) étant fixée, soit \(S_K\) l’ensemble fini des
canaux de la réalisation considérée. Ses éléments conservent les
données de route et de représentation qui distinguent encore les états :
feuillet, orientation, frontière, mémoire, saveur, couleur et composante de
spin lorsqu’ils appartiennent au secteur. Le quotient qui définit \(S_K\)
s’effectue avant de construire l’occupation ; il n’identifie pas deux canaux
par la seule égalité de leurs masses.

L’espace d’une particule est
\begin{equation}
 \mathcal H_K=\bigoplus_{s\in S_K}V_s,
 \qquad\langle e_{s,a},e_{t,b}\rangle=\delta_{st}\delta_{ab},
 \label{vi:vi:eq:espacio-canales-fock}
\end{equation}
où \(V_s\) porte la multiplicité de représentation déclarée.
Chaque \(e_{s,a}\) est un mode complet, non une cellule de l’atlas sans
étiquettes. Son produit intérieur positif appartient à cette réalisation
finie et explicite la normalisation utilisée ci-dessous.

La donnée de parité est un caractère compositionnel du transport,
\begin{equation}
 p:\operatorname{Mor}(\mathsf P_{\mathrm{TPK}})
 \longrightarrow\mathbb Z/2\mathbb Z,
 \qquad p(\gamma\eta)=p(\gamma)+p(\eta).
 \label{vi:vi:eq:caracter-paridad-fock}
\end{equation}
Sur les canaux où cette graduation est réalisée,
l’involution d’holonomie \(J_p\) satisfait
\(J_pe=(-1)^{p(e)}e\). Par conséquent,
\[
 \mathcal H_K=\mathcal H_{K,\bar0}\oplus\mathcal H_{K,\bar1},
 \qquad \Pi_{\bar0}=\frac{I+J_p}{2},\quad
 \Pi_{\bar1}=\frac{I-J_p}{2}.
\]
Les deux projecteurs sont orthogonaux car \(J_p^*=J_p\) et
\(J_p^2=I\). L’égalité \eqref{vi:vi:eq:caracter-paridad-fock} assure
que composer les transports conserve la graduation. Lorsqu’une réalisation
n’a pas fourni ce caractère, la complétion qui suit est une
construction relative à son choix explicite, non une attribution
statistique obtenue à partir de l’apparence d’une route.

\subsection{Échange gradué et projecteurs de permutation}
\label{vi:vi:subsec:intercambio-graduado}

Pour des vecteurs homogènes, on définit l’échange
\begin{equation}
 \tau_p(v\otimes w)=(-1)^{p(v)p(w)}w\otimes v.
 \label{vi:vi:eq:intercambio-graduado}
\end{equation}
Le signe est négatif exactement lorsque les deux modes appartiennent au
secteur impair. L’échange ne modifie pas les étiquettes d’un mode :
il transporte le mode complet d’une position tensorielle à une autre.

\Needspace{10\baselineskip}
\begin{proposition}[Représentation graduée des permutations]
\label{vi:vi:prop:permutaciones-graduadas}
Les échanges adjacents \(\tau_{p,j}\) sur
\(\mathcal H_K^{\otimes n}\) sont unitaires et satisfont
\[
\begin{gathered}
 \tau_{p,j}^2=I,\qquad
 \tau_{p,j}\tau_{p,k}=\tau_{p,k}\tau_{p,j}\quad(|j-k|>1),\\
 \tau_{p,j}\tau_{p,j+1}\tau_{p,j}
 =\tau_{p,j+1}\tau_{p,j}\tau_{p,j+1}.
\end{gathered}
\]
Ils définissent donc une représentation unitaire \(U_p\) de
\(\mathfrak S_n\).
\end{proposition}
\begin{proof}
Appliquer deux fois un échange donne le signe
\((-1)^{2p(v)p(w)}=1\) et restaure les deux vecteurs. Les opérateurs
qui agissent sur des paires disjointes commutent. Sur un triplet homogène
de degrés \(a,b,c\), les deux membres de la relation à trois
échanges produisent le même ordre final et le signe
\((-1)^{ab+ac+bc}\). Chaque opérateur permute une base orthonormale
en la multipliant par des signes de module un ; il est donc unitaire.
\end{proof}

La moyenne
\[
 P_n^p=\frac1{n!}\sum_{\pi\in\mathfrak S_n}U_p(\pi)
\]
est le projecteur orthogonal sur les tenseurs invariants par
l’échange gradué. Sur un espace homogène pair, il coïncide avec le
symétriseur ordinaire ; sur un espace homogène impair, il coïncide avec
l’antisymétriseur. Si \(U\) désigne la permutation sans signe,
\begin{equation}
 S_n=\frac1{n!}\sum_\pi U(\pi),\qquad
 A_n=\frac1{n!}\sum_\pi\operatorname{sgn}(\pi)U(\pi).
 \label{vi:vi:eq:simetrizadores-fock}
\end{equation}

\begin{proposition}[Symétrisation et antisymétrisation exactes]
\label{vi:vi:prop:proyectores-fock}
On a \(S_n^*=S_n^2=S_n\) et \(A_n^*=A_n^2=A_n\).
Pour \(n\ge2\), \(S_nA_n=0\). Leurs images sont, respectivement,
\(\operatorname{Sym}^n\mathcal H\) et \(\bigwedge^n\mathcal H\).
\end{proposition}
\begin{proof}
L’adjonction remplace chaque permutation par son inverse et laisse
les sommes invariantes. Dans le carré de chaque moyenne, toute
permutation \(\rho\) apparaît \(n!\) fois ; dans la seconde moyenne,
son coefficient est \(\operatorname{sgn}(\rho)\). Cela prouve
l’idempotence. Le produit des deux projecteurs inclut la somme
des signes de toutes les permutations, qui est nulle pour
\(n\ge2\). Enfin, leurs images sont les tenseurs invariants
et alternés par leurs propres formules de moyenne.
\end{proof}

\subsection{Algèbre universelle et normalisation des états d’occupation}
\label{vi:vi:subsec:fock-construccion}

Soit \(T(\mathcal H)=\bigoplus_{n\ge0}\mathcal H^{\otimes n}\)
l’algèbre tensorielle. La complétion algébrique graduée est le
quotient par l’idéal engendré par
\begin{equation}
 v\otimes w-(-1)^{p(v)p(w)}w\otimes v.
 \label{vi:vi:eq:ideal-fock}
\end{equation}
La propriété universelle du quotient signifie que toute application
linéaire graduée des canaux vers une algèbre qui vérifie ces
relations se prolonge de manière unique en un homomorphisme du
quotient. La règle d’échange, une fois déclarée, fixe ainsi le
produit des occupations sans recourir à des énergies cibles.

\begin{theorem}[Décomposition symétrique et extérieure]
\label{vi:vi:thm:fock-graduado}
Si \(d_0=\dim\mathcal H_{\bar0}\) et
\(d_1=\dim\mathcal H_{\bar1}\), la complétion graduée possède
la décomposition
\begin{equation}
 \mathcal F(\mathcal H)=
 \mathcal F_+(\mathcal H_{\bar0})\otimes
 \mathcal F_-(\mathcal H_{\bar1}),
 \label{vi:vi:eq:fock-producto}
\end{equation}
\[
 \mathcal F_+(\mathcal H_{\bar0})=
 \bigoplus_{m\ge0}\operatorname{Sym}^m\mathcal H_{\bar0},
 \qquad
 \mathcal F_-(\mathcal H_{\bar1})=
 \bigoplus_{k=0}^{d_1}\bigwedge^k\mathcal H_{\bar1}.
\]
Le secteur de nombre total \(n\) est
\[
 \mathcal F_n(\mathcal H)=
 \bigoplus_{k=0}^{\min(n,d_1)}
 \operatorname{Sym}^{n-k}\mathcal H_{\bar0}
 \otimes\bigwedge^k\mathcal H_{\bar1}.
\]
\end{theorem}
\begin{proof}
Des bases homogènes ordonnées étant fixées, les relations
\eqref{vi:vi:eq:ideal-fock} permettent de placer d’abord les facteurs pairs
puis les impairs. Les pairs commutent ; les impairs anticommutent,
et le carré d’un impair est nul car on travaille en caractéristique
zéro. Chaque monôme s’écrit comme
\[
 e_1^{m_1}\cdots e_{d_0}^{m_{d_0}}
 f_{i_1}\cdots f_{i_k},
 \qquad m_j\ge0,\quad i_1<\cdots<i_k.
\]
La réalisation par des tenseurs symétriques et alternés rend
ces monômes linéairement indépendants. Elle fournit l’application
inverse et démontre la décomposition. Les regrouper selon
\(m_1+\cdots+m_{d_0}+k=n\) donne le dernier énoncé.
\end{proof}

La complétion de Hilbert déclare orthonormaux les vecteurs
d’occupation normalisés \(|\boldsymbol m;I\rangle\), où
\(\boldsymbol m\in\mathbb N_0^{d_0}\) et
\(I\subset\{1,\ldots,d_1\}\). Dans le secteur bosonique de degré
\(n=\sum_jm_j\), la normalisation tensorielle est
\[
 |\boldsymbol m\rangle=
 \sqrt{\frac{n!}{\prod_jm_j!}}\,
 S_n(e_1^{\otimes m_1}\otimes\cdots\otimes
 e_{d_0}^{\otimes m_{d_0}}).
\]
Pour \(I=\{i_1<\cdots<i_k\}\), le vecteur extérieur est
\[
 |I\rangle=\frac1{\sqrt{k!}}\sum_{\pi\in\mathfrak S_k}
 \operatorname{sgn}(\pi)
 f_{i_{\pi(1)}}\otimes\cdots\otimes f_{i_{\pi(k)}}.
\]
Dans la première formule, il y a \(n!/\prod m_j!\) tenseurs distincts ;
dans la seconde, il y en a \(k!\). Les compter prouve que les deux normes valent un.
Le degré zéro est \(\mathbb C|0\rangle\) ; ce vide d'occupation
n'est identifié ni à la carte constitutive du vide électromagnétique
ni à une lacune du registre TPK.

\subsection{Création fermionique et dérivation des relations CAR}
\label{vi:vi:subsec:car-construccion}

Soit \(I\) un sous-ensemble ordonné de modes impairs et
\(r_j(I)=\#\{i\in I:i<j\}\). On définit
\begin{align}
 a_j^\dagger|I\rangle&=
 \begin{cases}
 0,&j\in I,\\
 (-1)^{r_j(I)}|I\cup\{j\}\rangle,&j\notin I,
 \end{cases}\notag\\
 a_j|I\rangle&=
 \begin{cases}
 (-1)^{r_j(I)}|I\setminus\{j\}\rangle,&j\in I,\\
 0,&j\notin I.
 \end{cases}
 \label{vi:vi:eq:creacion-fermi}
\end{align}
Le signe compte les transpositions nécessaires pour insérer ou retirer
le mode en maintenant l'ordre de la base extérieure. Il n'utilise ni sa masse,
ni sa charge, ni le signe d'une valeur scalaire publiée. Les matrices des
deux applications sont adjointes et leurs normes sont \(1\), sauf dans le
cas trivial sans modes impairs.

\begin{theorem}[Relations canoniques d’anticommutation]
\label{vi:vi:thm:car-construidas}
Les opérateurs de \eqref{vi:vi:eq:creacion-fermi} satisfont
\[
 \{a_i,a_j^\dagger\}=\delta_{ij}I,\qquad
 \{a_i,a_j\}=0,\qquad
 \{a_i^\dagger,a_j^\dagger\}=0.
\]
Ce sont des relations exactes sur la totalité de
\(\mathcal F_-(\mathcal H_{\bar1})\).
\end{theorem}
\begin{proof}
Pour \(i=j\), le produit \(a_j^\dagger a_j\) est l’identité
sur les vecteurs qui contiennent \(j\) et zéro sur les autres ;
\(a_ja_j^\dagger\) possède les actions complémentaires. Leur somme
est \(I\). Deux insertions ou deux retraits du même mode donnent zéro.
Pour \(i\ne j\), si une opération est empêchée, les deux termes
de l’anticommutateur correspondant sont nuls. Si les deux compositions
sont admissibles, elles produisent le même ensemble d’occupations. Changer
l’ordre modifie d’une unité le nombre de modes qui précèdent
l’une des insertions ou l’un des retraits ; les signes sont opposés. Cela
prouve les trois anticommutateurs sur chaque vecteur de la base.
\end{proof}

\begin{corollary}[Exclusion et occupation d’un mode complet]
\label{vi:vi:cor:pauli-ocupacion}
Pour \(N_j=a_j^\dagger a_j\),
\[
 N_j^*=N_j,\qquad N_j^2=N_j,\qquad
 \operatorname{Spec}(N_j)\subset\{0,1\}.
\]
Un même mode complet admet au plus une occupation fermionique.
\end{corollary}
\begin{proof}
Les CAR donnent
\[
 N_j^2=a_j^\dagger a_ja_j^\dagger a_j
 =a_j^\dagger(I-a_j^\dagger a_j)a_j=N_j,
\]
car \(a_j^2=(a_j^\dagger)^2=0\). L’autoadjonction est immédiate ;
une valeur propre d’un idempotent autoadjoint vaut \(0\) ou \(1\).
\end{proof}

La preuve construit CAR avant de déduire Pauli : elle n’introduit pas
l’exclusion comme restriction supplémentaire d’un tableau de particules.
Elle explique également les générateurs réels employés dans la
proposition~\ref{vi:vi:prop:car-real} : cette proposition agit maintenant
sur l’algèbre extérieure explicite. Son interprétation de Majorana
conserve les conditions de la section~\ref{vi:vi:sec:topologia}.

\subsection{Création bosonique, CCR et domaine des particules en nombre fini}
\label{vi:vi:subsec:ccr-construccion}

Sur la base bosonique normalisée, on définit
\begin{equation}
 b_j^\dagger|\boldsymbol m\rangle=
 \sqrt{m_j+1}\,|\boldsymbol m+\boldsymbol e_j\rangle,
 \qquad
 b_j|\boldsymbol m\rangle=
 \sqrt{m_j}\,|\boldsymbol m-\boldsymbol e_j\rangle,
 \label{vi:vi:eq:creacion-bose}
\end{equation}
avec une valeur nulle dans la seconde expression lorsque \(m_j=0\).
Les facteurs proviennent de la normalisation des tenseurs symétriques.
Le domaine commun \(\mathcal D_{\mathrm{fin}}\) est la somme directe
algébrique des secteurs de nombre : les vecteurs n’ont qu’un
nombre fini de composantes d’occupation non nulles. Il est dense et
invariant sous les applications précédentes.

\begin{theorem}[Relations canoniques de commutation]
\label{vi:vi:thm:ccr-construidas}
Dans \(\mathcal D_{\mathrm{fin}}\),
\[
 [b_i,b_j^\dagger]=\delta_{ij}I,\qquad
 [b_i,b_j]=[b_i^\dagger,b_j^\dagger]=0.
\]
Les opérateurs de nombre \(N_j^B=b_j^\dagger b_j\) satisfont
\(N_j^B|\boldsymbol m\rangle=m_j|\boldsymbol m\rangle\).
\end{theorem}
\begin{proof}
Pour un même mode,
\(b_jb_j^\dagger|\boldsymbol m\rangle=(m_j+1)|\boldsymbol m\rangle\)
et \(b_j^\dagger b_j|\boldsymbol m\rangle=m_j|\boldsymbol m\rangle\).
La différence est l’identité. Pour des modes distincts, les facteurs
affectent des coordonnées d’occupation différentes ; ils commutent donc.
Les formules de nombre découlent de l’application consécutive des deux
opérations de \eqref{vi:vi:eq:creacion-bose}.
\end{proof}

S’il existe un mode pair, \(b_j\) et \(b_j^\dagger\) sont des opérateurs
non bornés dans la complétion de Hilbert : leurs normes sur
\(|m_j\rangle\) croissent comme des racines de \(m_j\). Les égalités
précédentes ont donc un domaine explicite et ne sont pas présentées
comme des identités entre matrices finies sur tout l’espace bosonique.
Les opérateurs fermioniques agissent sur le second facteur de
\eqref{vi:vi:eq:fock-producto} ; les bosoniques, sur le premier. Sur
le domaine commun, les opérateurs de ces deux facteurs commutent.

\begin{proposition}[Défaut exact d’une coupure bosonique dure]
\label{vi:vi:prop:corte-bosonico}
Dans un mode, soit \(P_N\) la projection sur
\(|0\rangle,\ldots,|N\rangle\) et soit \(b^{(N)}=P_NbP_N\).
Sur cet espace fini,
\[
 [b^{(N)},(b^{(N)})^\dagger]
 =I-(N+1)|N\rangle\langle N|.
\]
\end{proposition}
\begin{proof}
Pour \(m<N\), la création et l’annihilation conservent le calcul
\((m+1)-m=1\). En \(|N\rangle\), la création tronquée est
nulle et le produit inverse vaut \(N\) ; le commutateur vaut
\(-N=1-(N+1)\). Ce sont tous les vecteurs de la base.
\end{proof}

La correction réside au dernier niveau de la coupure, non dans la loi
d’occupation. Un calcul fini peut vérifier les CCR loin de cette
frontière et doit conserver le terme de bord lorsqu’il inclut le
niveau terminal. Cette distinction est indispensable pour qu’un
certificat matriciel ne confonde pas une troncature avec la
complétion bosonique complète.

\subsection{Identité des modes complets et principe d'exclusion de Pauli}
\label{vi:vi:subsec:pauli-etiquetas-hojas}

L’exclusion se rapporte au même vecteur d’un mode complet. Deux
canaux peuvent partager une cellule, une masse ou une signature réduite et rester
linéairement indépendants par leur spin, leur feuillet, leur couleur,
leur saveur ou une autre coordonnée conservée par le transport. Ces deux
vecteurs admettent un produit extérieur non nul. Si un quotient
physique les identifie, ce quotient doit agir avant d’affirmer
combien de modes indépendants existent.

\begin{proposition}[Critère d’indépendance pour une occupation double]
\label{vi:vi:prop:pauli-gram}
Pour deux vecteurs d’un espace de canaux,
\[
 \|u\wedge v\|^2
 =\|u\|^2\|v\|^2-|\langle u,v\rangle|^2.
\]
En particulier, deux modes orthonormaux produisent un état extérieur
de norme un et deux vecteurs proportionnels produisent zéro.
\end{proposition}
\begin{proof}
La réalisation extérieure de degré deux est
\((u\otimes v-v\otimes u)/\sqrt2\). Développer sa norme donne
la formule. L’égalité dans Cauchy--Schwarz caractérise la
dépendance linéaire ; l’orthonormalité donne une norme au carré égale à un.
\end{proof}

Pour un canal orbital \(u\) et deux composantes de spin
orthogonales \(\xi_+,\xi_-\), les modes
\(u\otimes\xi_+\) et \(u\otimes\xi_-\) sont distincts et leur
produit extérieur peut être occupé. Il n’existe pas de double occupation
d’un même mode : il existe une occupation de deux composantes distinctes
de la représentation. Il en va de même pour deux feuillets lorsque ceux-ci
restent des composantes indépendantes de la réalisation.

Le mot électronique central peut être fixe sous \(C_M\), tandis que
la représentation de charge distingue électron et positron. Cette
coïncidence de mot n’identifie pas leurs modes complets ni
ne décide de leur occupation. L’involution de mot, l’inversion
cellulaire et l’échange des facteurs de Fock conservent les
domaines différenciés de la section~\ref{vi:vi:sec:conjugacion}.

Si \(r\) modes fermioniques orthonormaux possèdent la même valeur
d’une observable partielle, leur occupation totale peut prendre les valeurs
\(0,1,\ldots,r\). Chaque \(N_j\) reste un projecteur ;
la somme \(\sum_{j=1}^rN_j\) n’en est pas nécessairement un. Cette
conséquence distingue la dégénérescence spectrale et l’exclusion d’un
mode, deux propriétés qu’un chiffre de masse isolé ne sépare pas.

\subsection{Dimensions finies et transport par raffinement}
\label{vi:vi:subsec:fock-refinamiento}

\begin{proposition}[Recensement des occupations]
\label{vi:vi:prop:censo-fock}
Pour \(d\) modes,
\[
 \dim\bigwedge^n\mathbb C^d=\binom dn,
 \qquad
 \dim\operatorname{Sym}^n\mathbb C^d=\binom{d+n-1}{n}
 \quad(d\ge1).
\]
La première dimension est nulle pour \(n>d\). Dans l’espace
gradué, les dimensions de nombre total \(n\) s’obtiennent
en additionnant les produits des deux dimensions pour
\(k=0,\ldots,\min(n,d_1)\).
\end{proposition}
\begin{proof}
Une base extérieure se choisit au moyen de sous-ensembles de \(n\) modes
distincts. Une base symétrique se choisit au moyen de
\((m_1,\ldots,m_d)\in\mathbb N_0^d\) de somme \(n\).
Le placement de \(d-1\) séparateurs entre \(n\) occupations
dénombre ces compositions faibles et donne le second coefficient binomial.
La formule graduée découle de \eqref{vi:vi:eq:fock-producto} par
sommes directes et produits tensoriels.
\end{proof}

Soient \(T:\mathcal H_K\to\mathcal H_L\) les opérateurs d’entrelacement
de raffinement déclarés, qui conservent la graduation et les
données de canal correspondantes. Chacun induit
\begin{equation}
 \Gamma(T)=\bigoplus_{n\ge0}T^{\otimes n}|_{\mathcal F_n},
 \qquad T^{\otimes0}=I_{\mathbb C}.
 \label{vi:vi:eq:segunda-cuantizacion-transporte}
\end{equation}

\begin{theorem}[Compatibilité de la complétion avec le transport]
\label{vi:vi:thm:fock-naturalidad}
Si \(T\) conserve le degré, alors
\(T^{\otimes n}P_n^p=P_n^{p'}T^{\otimes n}\).
Si, en outre, \(\|T\|\le1\), \(\Gamma(T)\) est un opérateur
borné de norme au plus \(1\) sur Fock et
\[
 \Gamma(T_2T_1)=\Gamma(T_2)\Gamma(T_1),\qquad
 \Gamma(I)=I.
\]
Par conséquent, les carrés compatibles de raffinement
des canaux induisent des carrés compatibles d’occupations.
\end{theorem}
\begin{proof}
La conservation du degré fait commuter \(T^{\otimes n}\)
avec chaque échange adjacent ; par la
proposition~\ref{vi:vi:prop:permutaciones-graduadas}, il commute
avec les permutations et leur moyenne. Les restrictions
sont ainsi bien définies. Leur norme est bornée par
\(\|T\|^n\le1\), de sorte que la somme directe est bornée.
Les identités et les compositions se vérifient à chaque degré
et, par continuité, dans la complétion.
\end{proof}

La contractivité est une condition suffisante qui permet de traiter
la somme bosonique sans restrictions supplémentaires. Aux degrés finis,
les puissances existent pour toute application linéaire, mais leur
somme n’est pas toujours bornée : pour un mode bosonique et \(T=2I\),
\(\Gamma(T)|n\rangle=2^n|n\rangle\). Dans le secteur extérieur
fini, il n’y a que \(d_1+1\) degrés, de sorte que toute
application linéaire de ce secteur possède une seconde quantification
finie bornée. La différence relève du domaine, non de la causalité.

\begin{proposition}[Transport des créations et annihilations]
\label{vi:vi:prop:fock-transporte-adjunto}
Soit \(T:\mathcal H_K\to\mathcal H_L\) un transport borné
qui conserve le degré. Pour \(f\in\mathcal H_K\) et
\(g\in\mathcal H_L\), sur le domaine algébrique des particules
en nombre fini, on a
\begin{align}
 \Gamma(T)c_K^\dagger(f)
 &=c_L^\dagger(Tf)\Gamma(T),
 \label{vi:vi:eq:fock-transporte-creacion}\\
 c_L(g)\Gamma(T)
 &=\Gamma(T)c_K(T^*g).
 \label{vi:vi:eq:fock-transporte-aniquilacion}
\end{align}
Les indices indiquent la fibre sur laquelle agit chaque opérateur.
Si \(T\) est isométrique, la seconde identité se spécialise en
\[
 c_L(Tf)\Gamma(T)=\Gamma(T)c_K(f),
 \qquad
 N_{L,Tf}\Gamma(T)=\Gamma(T)N_{K,f},
\]
où \(N_{K,f}=c_K^\dagger(f)c_K(f)\) et, pour un mode unitaire,
\(N_{L,Tf}\) est son opérateur d’occupation transporté.
\end{proposition}
\begin{proof}
Dans un produit symétrique ou extérieur, appliquer \(T\) à chaque facteur
commute avec l’insertion de \(f\), qui devient \(Tf\). Cela
démontre \eqref{vi:vi:eq:fock-transporte-creacion}. L’annihilation
contracte chaque facteur avec le mode observé. Pour chaque facteur \(v\),
\[
 \langle g,Tv\rangle=\langle T^*g,v\rangle.
\]
Les signes extérieurs, ou gradués, sont conservés car \(T\)
préserve le degré ; les normalisations symétriques sont égales dans
les deux membres. La somme des contractions démontre
\eqref{vi:vi:eq:fock-transporte-aniquilacion}. Si \(T^*T=I\), on
substitue \(g=Tf\). La composition des identités de création
et d’annihilation donne ensuite celle d’occupation. Toutes les opérations
considérées préservent le domaine des particules en nombre fini ; lorsque
\(\|T\|\le1\), \(\Gamma(T)\) possède en outre l’extension bornée
du théorème~\ref{vi:vi:thm:fock-naturalidad}.
\end{proof}

L’identité générale conserve l’adjoint \(T^*\) et est valide sans
isométrie. La nécessité de la condition isométrique apparaît lors du
transport du même mode. Par exemple, si \(T=tI\), avec
\(0<t<1\), pour un mode unitaire \(f\) et le vide \(\Omega\),
\[
 c(tf)\Gamma(T)f=t^2\Omega,
 \qquad \Gamma(T)c(f)f=\Omega.
\]
La formule générale donne correctement
\(\Gamma(T)c(T^*tf)f=t^2\Omega\). On conserve ainsi le
transport par contractions, avec son adjoint explicite, et la
spécialisation isométrique de l’occupation.

\subsection{Symétrie d’échange et registre d’un composé}
\label{vi:vi:subsec:pauli-composicion}

La composition de la section~\ref{vi:vi:subsec:composicion-registro}
agit d’abord sur les histoires et leurs frontières. Fock organise
l’occupation des canaux ; \(\mathfrak C_{12}\) organise le
registre de liaison. Ce sont des opérations coordonnées qui ne se
remplacent pas l’une l’autre. En particulier, l’antisymétrisation ne
supprime pas la mémoire de l’ordre de couplage.

Soit \(\mathcal D\subset\mathcal H^{\otimes n}\) l’espace
engendré par les expressions de constituants admissibles pour une
interface. La linéarisation de la composition s’écrit
\(C_{\mathcal T}:\mathcal D\to\mathcal H_{\mathrm{comp}}\).
Lorsqu’un groupe fini \(G\) d’échanges conserve
l’admissibilité et la frontière, la compatibilité concrète est
\[
 C_{\mathcal T}U(g)=V(g)C_{\mathcal T},\qquad g\in G,
\]
où \(U,V\) sont des représentations unitaires et \(V\) transporte
les registres composés, y compris leurs données de feuillet et de mémoire.
L’équation est la condition
d’échange de cette réalisation ; elle ne se déduit pas du simple
recensement des expressions.

\begin{proposition}[Descente de la composition au secteur statistique]
\label{vi:vi:prop:composicion-estadistica}
Soit \(\chi:G\to\{+1,-1\}\) le caractère du secteur
d’échange. Sous la compatibilité précédente,
\[
 P_\chi^{\mathrm{comp}}C_{\mathcal T}
 =C_{\mathcal T}P_\chi^{\mathrm{in}},
\]
où
\[
 P_\chi^{\mathrm{in}}=\frac1{|G|}\sum_g\chi(g)U(g),
 \qquad
 P_\chi^{\mathrm{comp}}=\frac1{|G|}\sum_g\chi(g)V(g).
\]
\end{proposition}
\begin{proof}
Multiplier chaque égalité d’entrelacement par
\(\chi(g)/|G|\) et sommer donne exactement l’identité.
Les moyennes sont des projecteurs par le même calcul que dans
la proposition~\ref{vi:vi:prop:proyectores-fock}.
\end{proof}

Le résultat justifie la composition au sein d’un secteur symétrique
ou alterné lorsque la frontière possède le transport requis.
Si une permutation change le canal physique de liaison, on
conserve les arbres distincts et leurs interrelations ; on ne les
identifie pas par une symétrisation qui oublie cette frontière.

\subsubsection{Antisymétrie chromatique et reste de l’état baryonique}

Dans la réalisation baryonique, la
proposition~\ref{vi:vi:prop:singlete-barion} construit le vecteur
chromatique \(|\varepsilon\rangle\), qui change de signe
sous une transposition. Après regroupement des facteurs de l’état
complet, un vecteur factorisé a la forme
\[
 \Psi=|\varepsilon\rangle_{\mathrm{col}}
 \otimes\Phi_{\text{spin,saveur,orbital,feuillet}}.
\]
L’échange de deux constituants agit simultanément
sur les deux facteurs. Si \(U_{\mathrm{rest}}(\pi)\Phi=\Phi\),
alors l’état complet change de signe et appartient au
secteur fermionique. Réciproquement, pour un vecteur factorisé
avec le singulet chromatique non nul, l’alternance totale exige
que le reste soit symétrique sous ces mêmes permutations.
La preuve consiste à factoriser
\(U_{\mathrm{tot}}(\pi)=U_{\mathrm{col}}(\pi)\otimes
U_{\mathrm{rest}}(\pi)\) et à simplifier le signe chromatique.

Le mot « reste » inclut explicitement la saveur et le feuillet.
Échanger les couleurs sans transporter ces composantes n’est pas
la permutation de deux modes complets. Pour un état non
factorisable, on applique le projecteur alterné à l’espace
tensoriel complet, sans le remplacer par une règle scalaire.

\subsubsection{Compatibilité des projecteurs physiques}

Les projecteurs de couleur, de canal et de statistique peuvent être combinés
par produit lorsqu’ils commutent. Si \(P,Q\) sont des projecteurs
orthogonaux et \([P,Q]=0\), alors
\((PQ)^*=PQ\), \((PQ)^2=PQ\) et
\(\operatorname{im}(PQ)=\operatorname{im}P\cap\operatorname{im}Q\).
L’inclusion dans les deux images est directe ; pour un vecteur
dans l’intersection, \(PQv=v\). S’ils ne commutent pas, leur produit
ne représente généralement pas le projecteur des deux conditions
simultanées : la réalisation doit construire le sous-espace
commun ou l’ordre physique des opérations.

Cette précision conserve le sens des projecteurs
singulet et de Pauli qui interviennent dans les composés de la
section~\ref{vi:vi:sec:familias-compuestos}. Le couplage
sélectionne un secteur ; la masse composée est évaluée après
la construction du registre de liaison de ce secteur.

\subsection{Action de l’opérateur matériel et portée de l’identification physique}
\label{vi:vi:subsec:fock-alcance}

Soit \(B=B^*\) un opérateur d’un canal déjà obtenu à partir de la
dynamique interne et qui conserve la graduation. Au degré \(n\),
il agit par
\[
 d\Gamma_n(B)=\sum_{j=1}^n
 I^{\otimes(j-1)}\otimes B\otimes I^{\otimes(n-j)}.
\]
La somme commute avec les permutations graduées et se restreint
à \(\mathcal F_n\). Si \(B\ge0\), chaque restriction est positive.
Dans une base propre, sa valeur est la somme des occupations
multipliées par les valeurs propres de \(B\), avec les occupations
permises par les preuves précédentes. C’est l’action
d’une observable additive sur des canaux indépendants ;
elle n’identifie pas l’opérateur d’un composé lié à la
somme des masses isolées de ses constituants.

Dans la somme de Hilbert bosonique, \(d\Gamma(B)\) possède le
domaine naturel
\[
 \left\{\psi=(\psi_n):
 \sum_n\|d\Gamma_n(B)\psi_n\|^2<\infty\right\}.
\]
Chaque bloc est autoadjoint en dimension finie, et leur somme
directe sur ce domaine est autoadjointe. L’opérateur de liaison
ou de canaux ouverts n’est ajouté que lorsqu’il a été construit
par la dynamique correspondante. La formation d’un pôle
complexe suit le traitement de
\ref{vi:vi:subsec:resonancias-familias}.

On a donc obtenu des résultats algébriques complets :
projecteurs de symétrie, bases normalisées, CAR et CCR sur
leurs domaines, exclusion d’un mode, recensements d’occupation et
compatibilité avec les opérateurs d’entrelacement de raffinement. Le point
de départ de ces résultats est déclaré dans
\eqref{vi:vi:eq:espacio-canales-fock},
\eqref{vi:vi:eq:caracter-paridad-fock} et
\eqref{vi:vi:eq:intercambio-graduado}.

Lorsque la réalisation des rotations satisfait
\(U(2\pi)=(-1)^p\), la même graduation coordonne la
restauration de l’orientation à \(4\pi\) et la complétion
extérieure. Cette compatibilité n’est pas remplacée par une
coïncidence entre deux signes. L’identification avec le
théorème relativiste de spin--statistique nécessite en outre
la représentation lorentzienne, la localité et la positivité
de la réalisation physique. Les secteurs de tressage qui ne
passent pas aux permutations restent dans le codomaine de
la section~\ref{vi:vi:sec:topologia} ; ils ne sont pas forcés à entrer
dans cette complétion bosonique ou fermionique.
